% !TEX root = ../thesis.tex
% ==============================================================================
%  CHAPTER 1: INTRODUCTION (فصل اول: مقدمه) 
%             (chapters/chapter1_introduction.tex)
% ==============================================================================
%  This chapter introduces the research problem, problem formulation,
%  significance of the study, research objectives, novel contributions,
%  and the overall roadmap of the thesis.
%
%  TYPESETTING GUIDELINES:
%  - Equation refs: Use `\eqref{eq:...}` or `\ref{sec:...}`, both are adapted
%  - Chapter/Section/Figure/Table refs: Use `\ref{chap:...}`, `\ref{sec:...}`,
%                                            `\ref{fig:...}`, `\ref{tab:...}`
%  - Citations: Always prepend with a non-breaking space, e.g., `متن~\cite{ref}`
%  - English Footnotes: Use `\enfootnote{English Term}` for LTR footnotes
%  - Footnote Columns: Default is two-column (`\twocolumnfootnotes`). Switch to
%                      single-column via `\normalfootnotes` for long footnotes.
%  - Inline Code & File Paths: Use `\code{...}` for inline code, commands, and
%                              paths (ensures LTR monospace inside RTL text)
% ==============================================================================


\chapter{مقدمه}\label{chap:chapter1_introduction}

فصل اول پایان‌نامه به معرفی کلی موضوع پژوهش، بیان مسئله و انگیزه‌ی تحقیق، ضرورت انجام کار، اهداف، نوآوری‌ها و در نهایت ساختار کلی نوشتار اختصاص دارد. هدف اصلی این فصل، ارائه‌ی تصویری شفاف و منسجم از مسئله‌ی مورد مطالعه است تا خواننده پیش از ورود به مباحث تفصیلی در فصول بعد، با جایگاه و چارچوب کلی پژوهش آشنا شود. در نگارش این فصل، باید از ورود زودهنگام به جزئیات فنی یا بررسی تفصیلی پیشینه پرهیز کرد و تمرکز را بر تبیین روشن مبانی مسئله قرار داد.

متن این فصل علاوه بر تشریح استانداردهای نگارش قابلیت‌های پایه‌ای این قالب (شامل ارجاع به معادلات، بخش‌ها، مراجع علمی، کدهای درون‌متنی، پانوشت‌های دوزبانه و کادرها) را نشان می‌دهد.


\section{بیان مسئله و انگیزه پژوهش}\label{sec:1-1}

در بخش بیان مسئله، ابتدا زمینه‌ی موضوعی پژوهش ترسیم می‌شود و سپس چالش‌های موجود، محدودیت‌های روش‌های فعلی و شکاف علمی تبیین می‌گردد. ارجاع مستند به پژوهش‌های معتبر~\cite{sample_article, sample_book} برای نشان دادن اهمیت مسئله ضروری است. در این قالب، برای جلوگیری از شکستگی نامناسب شماره‌ی مرجع، همواره پیش از دستور ارجاع از فاصله‌ی نشکن (\lr{\textasciitilde}) استفاده می‌شود (مانند: متن\code{\textasciitilde{\textbackslash}cite\{\ldots\}}).

علاوه بر شرح کیفی، ارائه‌ی صورت‌بندی دقیق و ریاضی به شفافیت تعریف مسئله کمک شایانی می‌کند. برای نمونه، در یک مسئله‌ی بهینه‌سازی یا برآورد پارامتر، هدف می‌تواند یافتن بردار بهینه‌ی $\mathbf{x}^* \in \mathcal{X}$ در فضای پاسخ‌ها باشد، به‌گونه‌ای که تابع هزینه‌ی سامانه کمینه گردد:
\begin{equation}\label{eq:chap1_problem_formulation}
    \mathbf{x}^* = \arg\min_{\mathbf{x} \in \mathcal{X}} \left\{ f(\mathbf{x}) + \lambda \, g(\mathbf{x}) \right\}
\end{equation}
در رابطه‌ی~\eqref{eq:chap1_problem_formulation}، تابع $f(\mathbf{x})$ نشان‌دهنده‌ی خطای عملکرد سامانه، $g(\mathbf{x})$ تابع منظم‌ساز\enfootnote{regularizer} برای کنترل پیچیدگی یا ارتقای پایداری، و $\lambda > 0$ ضریب توازن میان این دو مؤلفه است.

\normalfootnotes{}
برای ارجاع به روابط از دستور \code{{\textbackslash}eqref\{\ldots\}}، برای ارجاع به بخش‌ها از دستور \code{{\textbackslash}ref\{\ldots\}} و برای درج واژه‌های تخصصی در پاورقی از \code{{\textbackslash}enfootnote\{\ldots\}} استفاده می‌شود. چیدمان پیش‌فرض قالب بر مبنای پانوشت‌های دو ستونی (\code{{\textbackslash}twocolumnfootnotes}) تنظیم شده است که برای واژه‌های کوتاه انگلیسی انتخابی بهینه است؛ اما چنانچه پانوشتی طولانی و تشریحی باشد یا ترجیح داده شود که پانوشت‌های یک صفحه به صورت تک‌ستونی تمام‌عرض درج گردند، می‌توان پیش از آن با دستور \code{{\textbackslash}normalfootnotes} حالت تک‌ستونی را فعال کرد و سپس با \code{{\textbackslash}twocolumnfootnotes} چیدمان دو ستونی را برای صفحات بعد بازگرداند.\footnote{برای نمونه، این پانوشت نشان می‌دهد که در صورت نیاز به نگارش توضیحی مفصل، چگونه می‌توان با فراخوانی \code{{\textbackslash}normalfootnotes} پانوشت را به صورت تک‌ستونی تمام‌عرض درآورد تا خوانایی متن حفظ شده و از فشردگی نامناسب در ستون‌های باریک جلوگیری گردد.} همچنین برای درج کدهای درون‌متنی و اسامی فایل‌ها از دستور \code{{\textbackslash}code\{\ldots\}} استفاده می‌گردد؛ این دستور با تلفیق قلم تک‌فاصله (\code{{\textbackslash}texttt\{\}}) و جهت چپ‌به‌راست، از به‌هم‌ریختگی علائم نگارشی در متن راست‌به‌چپ فارسی جلوگیری می‌کند.
\twocolumnfootnotes{}


\section{اهمیت و ضرورت انجام پژوهش}\label{sec:1-2}

در این بخش، پژوهشگر به این پرسش پاسخ می‌دهد که چرا پرداختن به موضوع پژوهش دارای اهمیت علمی و کاربردی است. توجیه ضرورت تحقیق باید مبتنی بر استدلال‌های منطقی، شواهد علمی و تبیین نیازهای عملیاتی باشد. ضرورت یک پژوهش می‌تواند در چند محور خلاصه شود. برای مثال، برطرف ساختن ناکارآمدی یا محدودیت‌های روش‌های پیشین، کاهش پیچیدگی‌های محاسباتی یا هزینه‌های اجرایی برای مقیاس‌پذیری سامانه‌ها و یا دستیابی به پاسخ‌هایی دقیق‌تر و تعمیم‌پذیرتر. ارائه‌ی این موارد در قالب بندهای پیوسته و مستدل، چارچوب منطقی مناسبی برای هدایت خواننده به سوی فصول بعدی فراهم می‌سازد.


\section{اهداف و پرسش‌های پژوهش}\label{sec:1-3}

هدف اصلی پژوهش باید به‌صورت صریح، شفاف و سنجش‌پذیر بیان شود تا جهت‌گیری کل پایان‌نامه مشخص باشد. در کنار هدف اصلی، طرح پرسش‌های کلیدی به تعیین دامنه‌ی دقیق تحقیق کمک می‌کند.

برای نمونه، هدف این پایان‌نامه می‌تواند ارائه‌ی الگوریتمی کارآمد برای حل مسئله‌ی بخش~\ref{sec:1-1} باشد. در راستای این هدف، پرسش‌های محوری تحقیق عبارتند از: چگونگی کاهش پیچیدگی محاسباتی بدون افت کارایی، چگونگی تضمین پایداری، و نحوه‌ی ارزیابی عملکرد روش در مقایسه با روش‌های استاندارد.


\section{نوآوری‌ها و مشارکت‌های علمی}\label{sec:1-4}

در این بخش، دستاوردهای اصیل و سهم علمی مشخص پایان‌نامه به‌روشنی فهرست می‌شوند تا مرز میان کارهای گذشته و نوآوری نویسنده برای داوران و خوانندگان شفاف باشد. در این قالب، برای جلب توجه مخاطب به نتایج کلیدی، می‌توان از کادر آکادمیک (\code{{\textbackslash}academicbox}) استفاده کرد که برای نخستین بار در زیر به نمایش درآمده است:

\vspace{0.3cm}
\academicbox{مشارکت‌های کلیدی پژوهش}{%
دستاوردهای اصلی این پایان‌نامه در محورهای زیر خلاصه می‌شود:\vspace{-0.4cm}
\begin{itemize}
    \item ارائه‌ی چارچوب تحلیلی یکپارچه برای صورت‌بندی مسئله در شرایط ابعاد بالا؛\vspace{-0.4cm}
    \item طراحی روشی با پیچیدگی محاسباتی کمتر در مقایسه با روش‌های شاخص موجود؛\vspace{-0.4cm}
    \item ارزیابی جامع عملکرد روش پیشنهادی بر روی معیارهای استاندارد و اعتبارسنجی تجربی.
\end{itemize}
}
\vspace{0.3cm}


\section{ساختار کلی پایان‌نامه و راهنمای پیوست‌ها}\label{sec:1-5}

بخش پایانی فصل اول، نقشه‌ی راه خواننده برای مطالعه‌ی فصول و پیوست‌های سند است. نوشتار حاضر بر اساس روندی از مبانی پایه تا جمع‌بندی نتایج تنظیم شده است. \textbf{فصل دوم (\nameref{chap:chapter2_preliminaries})} به ارائه‌ی مفاهیم پایه، نمادگذاری‌های ریاضی، محیط‌های رسمی تعریف و قضیه، و نمودارهای برداری اختصاص دارد. \textbf{فصل سوم (\nameref{chap:chapter3_literature})} اصول مرور ساختاریافته‌ی ادبیات پژوهش، نقد مقایسه‌ای رویکردها و تبیین شکاف تحقیقاتی را پوشش می‌دهد. \textbf{فصل چهارم (\nameref{chap:chapter4_methodology})} چارچوب پیشنهادی، فرمول‌بندی روابط و محیط نگارش الگوریتم را تشریح می‌کند. \textbf{فصل پنجم (\nameref{chap:chapter5_experiments})} استانداردهای ارائه‌ی نتایج تجربی، آزمایش‌های مقایسه‌ای و تحلیل‌های عملکردی را ارائه می‌نماید. \textbf{فصل ششم (\nameref{chap:chapter6_conclusions})} به جمع‌بندی یافته‌ها، بیان واقع‌بینانه‌ی محدودیت‌ها و پیشنهاد مسیرهای پژوهشی آینده می‌پردازد.

علاوه بر فصول فوق، نکات فنی، مدیریتی و ساختاری قالب در پیوست الف (\ref{chap:appendix}) مستند شده‌اند. در بخش~\ref{sec:app_title_and_forms}، نحوه‌ی تنظیم اطلاعات فراداده در فایل \code{frontmatter/info.tex}، فواصل صفحات عنوان و سازوکار صدور و ادغام فرم‌های اداری تصویب‌نامه و اظهارنامه تشریح شده است. در بخش~\ref{sec:app_dictionary_and_glossary}، معماری واژه‌نامه‌ی دوزبانه در فایل \code{backmatter/words.tex} و خط‌لوله‌ی مرتب‌سازی خودکار الفبای فارسی با اسکریپت پایتون مستند گردیده است. در بخش~\ref{sec:app_bibliography_rules}، ساختار و قواعد سبک کتاب‌شناسی \code{unsrtabbrv-modern.bst}، پیوندهای برخط (\lr{DOI / URL}) و انواع مدخل‌های مراجع بیان شده است.
